\begin{abstract}

Generative Adversarial Networks (GANs) have emerged as a powerful paradigm in generative modelling, capable of producing high-fidelity data samples across various domains. However, their performance remains tightly coupled with the quality of latent representations, which are traditionally derived from classical noise. In this study, we explore the frontier of quantum-enhanced generative modelling through Hybrid Quantum-Classical GANs (HQCGANs), where quantum circuits are used to generate latent vectors, while a classical neural network functions as the discriminator. The central hypothesis posits that quantum latent spaces, even with a few qubits, may introduce diversity and richness that classical priors lack, thereby enhancing generative fidelity.

We implement and evaluate a classical GAN alongside three HQCGAN variants using 3, 5, and 7 qubits, respectively. The quantum circuits are simulated via Qiskit's AerSimulator and incorporate realistic noise channels-depolarising, amplitude damping, and readout noise—reflecting the conditions of near-term quantum devices. Binary MNIST (digits 0 and 1) was chosen as the target dataset to align with the low-dimensional latent representations imposed by current quantum hardware limitations. Standard preprocessing and balanced training sets were used to ensure parity between models.

Each model was trained for 150 epochs using the Adam optimiser and evaluated with industry-standard metrics: Fréchet Inception Distance (FID) and Kernel Inception Distance (KID). Training dynamics and sample quality progression were analysed through visualisations. The results indicate that while the classical GAN achieved the best overall FID and KID scores, the HQCGANs, particularly the 7-qubit variant, obtained competitive performance, demonstrating the potential of quantum priors to match their classical counterparts in certain configurations. Notably, the 3-qubit model plateaued earlier, underscoring the representational limitations of smaller quantum latent spaces.

Efficiency was also evaluated in terms of epoch-level training time and quantum shot usage. Despite requiring additional resources for latent sampling, HQCGANs were found to be computationally feasible, with total shot counts kept constant and training times only moderately affected.

This study validates the viability of using noisy quantum circuits as latent priors within GAN architectures, and provides empirical support for theoretical claims regarding the expressive capacity of quantum-generated distributions. Limitations include reliance on simulated quantum environments, low qubit counts, and the absence of comparisons with more sophisticated GAN variants. Future work should consider real hardware implementations, complete quantum discriminators, and integration with quantum error mitigation techniques.

In conclusion, HQCGANs offer a promising hybrid approach that capitalises on quantum randomness and classical learning, potentially paving the way for more efficient and expressive generative models in the NISQ era. This work serves as a stepping stone towards understanding how quantum computing can be integrated meaningfully into modern deep learning workflows.
\end{abstract}
